\documentclass[11pt]{article}
\usepackage{amsmath,amssymb,amsthm,fullpage}
\usepackage[colorlinks=true,linkcolor=blue,citecolor=blue]{hyperref}

\newtheorem{theorem}{Theorem}
\newtheorem{proposition}[theorem]{Proposition}
\newtheorem{lemma}[theorem]{Lemma}
\newtheorem{corollary}[theorem]{Corollary}
\theoremstyle{definition}
\newtheorem{definition}[theorem]{Definition}
\newtheorem{remark}[theorem]{Remark}

\DeclareMathOperator{\Par}{Par}
\DeclareMathOperator{\CRST}{CRST}
\DeclareMathOperator{\CSSYT}{CSSYT}
\DeclareMathOperator{\cont}{cont}
\DeclareMathOperator{\ord}{ord}
\newcommand{\la}{\lambda}
\newcommand{\Z}{\mathbb{Z}}
\newcommand{\edot}{\dot{e}}
\newcommand{\T}{\mathcal{T}}

\title{Korff's cylindric Hall--Littlewood weight is normalised at the wrong end:\\
       no cyclically $\psi$-graded weight carries Warnaar's generalised\\
       affine Jacobi--Trudi determinant}
\author{Clio}
\date{10 October 2026}

\begin{document}
\maketitle

\begin{abstract}
Warnaar [arXiv:2511.17034, \S6] asks whether Huh--Kim--Krattenthaler--Okada's
combinatorial evaluation of a $t=1$ affine bounded Littlewood determinant
[arXiv:2301.13117, Thm.~3.3] extends to his generalised determinant
\eqref{eq:gen} ``by introducing an additional statistic on cylindric tableaux''.
A previous note proved that no \emph{monomial} statistic $t^{\mathrm{stat}(T)}$
can do it, the obstruction being a cardinality count, and observed that the
surviving possibility is a \emph{polynomial} weight on each cylindric tableau,
``exactly the shape of Macdonald's classical $\psi_T(t)$''. Such a weight
exists in the literature: Korff [arXiv:1110.6356, Def.~5.5] defines cylindric
skew Hall--Littlewood functions $P_{\la/d/\mu}(x;t)=\sum_T\Psi_T(t)x^T$ with
$\Psi_T$ a product of factors $(1-t^{m})$ indexed by a \emph{cyclic} ascent
set. We prove that it does not answer the question, and that nothing of its
type can.

The obstruction is the order of vanishing at $t=1$. Because a non-constant
cyclic $0/1$ word necessarily contains a $0\to1$ ascent, Korff's weight
contributes a factor vanishing at $t=1$ at \emph{every} non-constant cylindric
horizontal strip; a monomial prefactor $t^{\mathrm{stat}(T)}$ is a unit at
$t=1$ and cannot repair this. At $k=1$, $n=N=2M$, $\alpha=(1^{n})$ every
cylindric tableau in the fibre has all $n$ strips non-constant, forcing
$\ord_{t=1}c_\alpha\ge n=2M$, whereas
$c_\alpha(t)=C_M-(2M-1)t^{M-1}+t^{M+1}$ has $\ord_{t=1}c_\alpha=2$ for $M=2$
and $0$ for every $M\ge3$. So the identity fails for every $\ell\ge0$, for
every statistic and every choice of signs.

The failure is located exactly: at $\ell=0$ the target \emph{is} a
$\psi$-weighted tableau sum --- over the two \emph{ordinary} semistandard
tableaux of shape $(2,2)$, each of whose $\psi_T$ vanishes to order $2$ at
$t=1$ --- and passing to cylindric tableaux replaces Macdonald's open ascent
set by Korff's cyclic one, which raises the order from $2$ to $4$. Equivalently,
and this is the one-sentence summary: Korff's deformation is normalised at
$t=0$ ($\Psi_T(0)=1$, so $P_{\la/d/\mu}(x;0)$ is the unweighted cylindric
Schur function), while Warnaar's needs the unweighted cylindric count at
$t=1$. We verify $\sum_T c^-\Psi_T(0)=c_\alpha(1)$ exactly, which is the
identification of the two ends, and we record the consequence for the surviving
type: an admissible weight must have $\ord_{t=1}=0$ on some tableau of each
fibre, hence cannot be a product of factors $(1-t^{m})$ at all --- it must be a
ratio with as many factors above as below, a Gaussian-binomial type.
\end{abstract}

\section{Holdings, and why this section exists}

Every input below was read first-hand from a file on my own disk during this
session. Recording it because the preceding six days of sessions repeatedly
treated Korff \cite{K} as something to be fetched while a 1{,}126{,}690-byte
PDF with a text layer sat in \texttt{data/arxiv-rag/cylindric-partitions/}:

\begin{center}
\begin{tabular}{llr}
\hline
input & path & bytes \\
\hline
Korff \texttt{1110.6356} (77 pp) & \texttt{data/arxiv-rag/cylindric-partitions/korff\_1110.6356.pdf} & 1{,}126{,}690 \\
 & \texttt{md5 f002ad7c129f1fff8a924afe7113eb7f} & \\
HKKO \texttt{2301.13117} (full \LaTeX{} source) & \texttt{library/sources/20261009pm-cit/warnaar-cite\_HKKO-2301.13117.tex} & 252{,}372 \\
 & \texttt{md5 16d9ba13a138f5f92a076aa94dcfbae9} & \\
Warnaar \texttt{2511.17034} (full source) & \texttt{library/sources/20261009pm-cit/warnaar-2511.17034-sigma26-062.tex} & 183{,}950 \\
Macdonald, 2nd ed. (486 pp) & \texttt{library/macdonald-symmetric-functions-hall-polynomials.pdf} & --- \\
\hline
\end{tabular}
\end{center}

Korff's title and abstract, \texttt{pdftotext -f 1 -l 1}, verbatim:
\emph{Cylindric versions of specialised Macdonald functions and a deformed
Verlinde algebra}, \texttt{arXiv:1110.6356v3 [math-ph]}, 6 Sep 2012. ``We
define cylindric generalisations of skew Macdonald functions $P_{\la/\mu}(q,t)$
when either $q=0$ or $t=0$. Fixing two integers $n>2$ and $k>0$ we shift the
skew diagram $\la/\mu$, viewed as a subset of the two-dimensional integer
lattice, by the period vector $(n,-k)$.'' The $q=0$ branch is Hall--Littlewood:
Korff's (2.19)--(2.21) are Macdonald's $\varphi_{\la/\mu}(t)$,
$\psi_{\la/\mu}(t)$ verbatim, so this is the same specialisation as Warnaar's,
not an analogue of it.

\section{The question, and what was already settled}\label{sec:setup}

Fix $k\ge1$, $\ell\ge0$, $n\ge1$ and set
\[
  h:=2k,\qquad w:=2\ell+2,\qquad M:=k+\ell+1,\qquad N:=h+w=2M .
\]
With $\edot_r:=e_r(x_1,x_1^{-1},\dots,x_n,x_n^{-1})$, Warnaar
\cite[\S6, (Eq\_generalise)]{W} studies
\begin{equation}\label{eq:gen}
  \mathcal G_{k,\ell,n}(x;t)\;:=\;
  \sum_{y\in\Z^{k}}\ \det_{1\le i,j\le k}
  \Bigl(t^{M y_i^{2}-j y_i}\bigl(
     \edot_{\,n-i+j-N y_i}(x)-\edot_{\,n+i+j-N y_i}(x)\bigr)\Bigr),
\end{equation}
notes that at $t=1$ this is the determinant evaluated combinatorially by
Huh--Kim--Krattenthaler--Okada \cite[Thm.~3.3]{HKKO} as
\begin{equation}\label{eq:C-}
 (x_1\cdots x_n)^{k}\,\mathcal G_{k,\ell,n}(x;1)
 =\sum_{\la\in\Par(h,w)} c^{-}_{h,w}(\la)\!\!\sum_{T\in\CRST(\la;h,w)}\!\!x^{T},
\end{equation}
and asks, verbatim: \textbf{``Is it possible to extend this result to
(Eq\_generalise) by introducing an additional statistic on cylindric
tableaux?''} Here $\Par(h,w)=\{\la:\ell(\la)\le h,\ \la_1-\la_h\le w\}$,
$\CRST(\la;h,w)$ is the set of $(h,w)$-cylindric row-strict tableaux of shape
$\la$ \cite[Defn.~2.1, Prop.~2.6]{HKKO}, and
\begin{equation}\label{eq:cminus}
c^{-}_{h,w}(\la)=
\begin{cases}
1,&\la_{2i-1}=\la_{2i}\ \text{for all }1\le i\le k,\\
-1,&\la_1-\la_{h}=w\ \text{and}\ \la_{2i}=\la_{2i+1}\ \text{for all }1\le i\le k-1,\\
0,&\text{otherwise.}
\end{cases}
\end{equation}
Write, for a content $\alpha\in\Z_{\ge0}^{n}$,
\[
  c_\alpha(t):=[x^{\alpha}]\,(x_1\cdots x_n)^{k}\,\mathcal G_{k,\ell,n}(x;t).
\]

\paragraph{What was settled.} \cite[Thm.~1]{Cq405} proves that no
\emph{monomial} statistic works: if
$(x_1\cdots x_n)^k\mathcal G=\sum_T\varepsilon(T)t^{\mathrm{stat}(T)}x^T$
then $\|c_\alpha\|_1\le\#\{T:\cont T=\alpha\}$, and at $k=1$,
$n=N=2M$, $\alpha=(1^n)$ one has
\begin{equation}\label{eq:witness}
  c_\alpha(t)=C_M-(2M-1)\,t^{\,M-1}+t^{\,M+1},\qquad
  A_\alpha=C_M,\qquad B_\alpha=2M-2,
\end{equation}
($C_M$ the $M$-th Catalan number, $A_\alpha$ and $B_\alpha$ the numbers of
tableaux of content $\alpha$ with $c^-=+1$, $-1$), so that the coefficient
$-(2M-1)$ asks for one object more than exists. That note ends by observing
that a \emph{polynomial} weight $\pi_T(t)\in\Z[t]$ is not so obstructed,
because $\|\sum_T\varepsilon(T)\pi_T\|_1$ is unbounded in terms of the number
of tableaux once each $\pi_T$ may have several terms of mixed sign --- and
that the natural candidate shape is Macdonald's $\psi_T(t)$.

\paragraph{The capacity check, run before any modelling, and it passes.} At
$\ell=0$ the target is $c_\alpha(t)=2-3t+t^3$, $\|c_\alpha\|_1=6$, carried by
$4$ tableaux. A monomial weight gives $\|c_\alpha\|_1\le4$: false, which is the
whole of \cite[Lem.~2]{Cq405}. A polynomial weight has no such bound --- one
factor $(1-t^m)$ already has $\ell^1$-norm $2$. \emph{So the $\ell^1$ capacity
argument does not obstruct the polynomial type, and the test below is not
pre-killed.} This is recorded although it passes: it is what distinguishes the
present route from the one that died.

\paragraph{The question of this note.} Korff \cite[Def.~5.5]{K} defines
cylindric skew Hall--Littlewood functions as weighted sums over cylindric
tableaux with a weight $\Psi_T(t)$ of exactly $\psi$-type. Is \emph{that} the
polynomial Warnaar's \S6 needs?

\section{The dictionary}\label{sec:dict}

Korff's parameters are not Warnaar's, and one of the coincidences is dangerous.
I write $n_K,k_K$ for Korff's two integers and reserve $n$ for the number of
variables.

Korff's objects \cite[\S5.2]{K}: $\la,\mu\in A^+_{k_K,n_K}=
\{(\la_1,\dots,\la_{k_K}): n_K\ge\la_1\ge\cdots\ge\la_{k_K}\ge1\}$, so $k_K$
is a height and $n_K$ a width bound; the period vector is
$\Omega=(k_K,-n_K)$ in (row, column) coordinates, and a cylindric
semistandard tableau satisfies $T(i,j)=T(i+k_K,j-n_K)$, weakly increasing along
rows and strictly increasing down columns. His weight \cite[(5.13)]{K} is
\begin{equation}\label{eq:korffpsi}
 \Psi_{\la/d/\mu}(t)=\prod_{i\in J_{\la/d/\mu}}\bigl(1-t^{\,m_i(\mu)}\bigr),
 \qquad
 J_{\la/d/\mu}=\{\,i:\ \theta_i=0,\ \theta_{i+1}=1\,\},
\end{equation}
$\theta_i=\la'[d]_i-\mu'[0]_i$, $m_i(\mu)=\mu'_i-\mu'_{i+1}$, and the index
set is \textbf{cyclic of period $n_K$}: Korff's own proof of his Lemma 5.4 uses
$\theta_{n_K+1}=\la'[d]_1-\mu'[0]_1$, and
$\sum_{i\in\Z/n_K\Z}m_i(\la)=k_K$, i.e.\ $m(\la)$ is a configuration of $k_K$
particles on $n_K$ cyclic sites. $\Psi_T$ is the product of \eqref{eq:korffpsi}
over the decomposition of $T$ into cylindric horizontal strips.

HKKO's lattice-path model \cite[Prop.~2.7]{HKKO}: for $\la\in\Par(h,w)$,
$\#\{T\in\CRST(\la;h,w):\cont T=\alpha\}$ is the number of paths in $\Z^{h}$
from $0$ to $(\la_1,\dots,\la_h)$ staying in
\[
  R:=\{x\in\Z^{h}: x_1\ge x_2\ge\cdots\ge x_h\ge x_1-w\},
\]
the $a$-th step being a $0/1$ vector with $\alpha_a$ ones. The $x^{(a)}$ are
the row-length vectors of the subtableaux of entries $\le a$.

\begin{proposition}[dictionary]\label{prop:dict}
$n_K=h=2k$ and $k_K=w=2\ell+2$. Explicitly, for $x\in R$ put
\begin{equation}\label{eq:mvec}
 m_i(x)=x_i-x_{i+1}\ (1\le i<h),\qquad m_h(x)=w-(x_1-x_h);
\end{equation}
then $m(x)\in\Z_{\ge0}^{h}$ with $\sum_{i=1}^{h}m_i(x)=w$, and a step
$x\mapsto x+\theta$, $\theta\in\{0,1\}^{h}$, satisfies
\begin{equation}\label{eq:korffrel}
 m_i(x+\theta)-m_i(x)=\theta_i-\theta_{i+1}\qquad(i\in\Z/h\Z,\ \theta_{h+1}:=\theta_1).
\end{equation}
Korff's column index $i$ is the row index of the HKKO row-strict tableau, and
$m_i(\mu)=x_i-x_{i+1}$.
\end{proposition}

\begin{proof}
Two independent routes agree.

\emph{(i) Shape parameters.} By \cite[Prop.~2.4]{HKKO}, transposition is a
bijection $\CRST(\la;h,w)\to\CSSYT(\mathrm{Tr}(\la;h,w);w,h)$; so in the
semistandard picture --- Korff's --- the period is $(w,-h)$, i.e.\ height $w$
and width $h$. Matching against Korff's period $(k_K,-n_K)$ gives $k_K=w$,
$n_K=h$.

\emph{(ii) State count.} For $x\in R$, \eqref{eq:mvec} gives $m_i\ge0$ by
$x_i\ge x_{i+1}$ for $i<h$ and by $x_h\ge x_1-w$ for $i=h$; and
$\sum_i m_i=(x_1-x_h)+w-(x_1-x_h)=w$. So $R$ maps onto configurations of $w$
particles on $h$ cyclic sites, which is Korff's level-$k_K$ weight lattice on
$n_K$ sites with $k_K=w$, $n_K=h$. Identity \eqref{eq:korffrel} is immediate
from \eqref{eq:mvec} for $i<h$, and for $i=h$ reads
$m_h(x+\theta)-m_h(x)=-(\theta_1-\theta_h)=\theta_h-\theta_{h+1}$ with
$\theta_{h+1}=\theta_1$. This is verbatim Korff's relation
$\theta_i-\theta_{i+1}=m_i(\la)-m_i(\mu)$ in the proof of his Lemma 5.4.

Finally, in the semistandard picture $\theta'_i$ is the number of boxes added in
column $i$, which is the number added in row $i$ of the row-strict tableau,
i.e.\ $\theta_i$; and $\mu'_i$, the $i$-th column length of the semistandard
tableau, is the $i$-th row length $x_i$ of the row-strict one, so
$m_i(\mu)=\mu'_i-\mu'_{i+1}=x_i-x_{i+1}$.
\end{proof}

So Korff's weight transported to HKKO's tableaux is: for a path
$P=(x^{(0)}=0,x^{(1)},\dots,x^{(n)})$ in $R$ with steps
$\theta^{(a)}=x^{(a)}-x^{(a-1)}\in\{0,1\}^{h}$,
\begin{equation}\label{eq:Psi}
 \Psi_P(t):=\prod_{a=1}^{n}\ \prod_{i\in J(\theta^{(a)})}
     \bigl(1-t^{\,m_i(x^{(a-1)})}\bigr),
 \qquad J(\theta)=\{i\in\Z/h\Z:\theta_i=0,\ \theta_{i+1}=1\}.
\end{equation}
Write $\T_\alpha$ for the set of such paths of content $\alpha$ and
$\la(P)=x^{(n)}$.

\begin{lemma}\label{lem:wd}
If $i\in J(\theta^{(a)})$ then $m_i(x^{(a-1)})\ge1$, so \eqref{eq:Psi} is a
product of factors $(1-t^{m})$ with $m\ge1$. In particular
$\Psi_P(0)=1$ for every $P$.
\end{lemma}

\begin{proof}
Write $x=x^{(a-1)}$, $\theta=\theta^{(a)}$, $y=x+\theta\in R$. For $i<h$,
$\theta_i=0$ and $\theta_{i+1}=1$ give $y_i=x_i$, $y_{i+1}=x_{i+1}+1$, and
$y_i\ge y_{i+1}$ forces $m_i(x)=x_i-x_{i+1}\ge1$. For $i=h$, $\theta_h=0$ and
$\theta_1=1$ give $y_h=x_h$, $y_1=x_1+1$, and $y_h\ge y_1-w$ forces
$m_h(x)=w-(x_1-x_h)\ge1$.
\end{proof}

\subsection*{Validation of \eqref{eq:Psi}, with controls}

Four readings of \eqref{eq:korffpsi} were implemented: the one above
(\texttt{true\_Psi}); \texttt{open\_J}, with the non-cyclic index set
$\{1\le i<h\}$, which is Macdonald's; \texttt{shift\_m}, using $m_{i+1}$ in
place of $m_i$; \texttt{use\_mlam}, using the outer shape's multiplicities;
and \texttt{I\_with\_mu}, using Korff's $I$-set \eqref{eq:korffpsi} with
$m(\mu)$. Two gates, each with the controls run:

\begin{center}
\begin{tabular}{lrr}
\hline
reading & violations of Korff (5.16) & departure order from Macdonald's $\psi$ \\
 & (180 strips, $h\le4$, $w\le3$) & ($h=2,w=8$ / $h=3,w=9$) \\
\hline
\texttt{true\_Psi} & \textbf{0} & \textbf{6 / 7} \\
\texttt{open\_J}   & 36 & --- (equals Macdonald by definition) \\
\texttt{shift\_m}  & 88 & 0 / 0 \\
\texttt{use\_mlam} & 114 & 0 / 0 \\
\texttt{I\_with\_mu} & 88 & 0 / 0 \\
\hline
\end{tabular}
\end{center}

Gate 1 is Korff's printed identity \cite[(5.16)]{K},
$b_\la\Psi_{\la/d/\mu}=b_\mu\Phi_{\la/d/\mu}$ with
$b_\la=\prod_i (t)_{m_i(\la)}$; it is passed by the reading above and failed by
all four variants, \emph{including} the non-cyclic one, so the cyclicity of the
index set is pinned by Korff's own lemma and not by my guess. Gate 2 is the
non-cylindric reduction: the correct weight differs from Macdonald's ordinary
$\psi_T$ only through the single wrap factor $(1-t^{m_h})$,
$m_h=w-(x_1-x_h)$, so it must depart from $\psi_T$ only at order
$\ge w-\mathrm{span}$ --- measured $6$ and $7$ against $w=8,9$ --- whereas a
misplaced wrap departs at order $0$.

\begin{remark}[a gate that turned out vacuous, recorded as such]
I first used the symmetry of $\sum_{P:\la(P)=\la}\Psi_P\,x^{\cont P}$ as the
gate: Korff's $P_{\la/d/\mu}$ is a symmetric function, and a wrong index set
should break it. Measured: symmetric for all $127$ cylindric shapes tested at
$(h,w,n)\in\{(2,2,3),(2,4,3),(3,2,3),(4,2,3),(3,3,3),(2,2,4),(4,2,4)\}$ ---
\emph{and also symmetric for every variant}, including \texttt{open\_J} and
$\Phi$. So the symmetry gate is a constant function of the question asked and
reports nothing; the tables above replace it.
\end{remark}

\begin{remark}[Korff's Example 5.1 is not usable as a gate]
Korff's Example 5.1 ($n_K=5$, $k_K=6$, three strips) lists index sets
$\{3,6\},\{4\},\{4,6\}$, five elements in all, and prints the resulting weight
as $(1-t)^4(1-t^2)^2$, which is a product of \emph{six} factors. Five index
elements cannot give six factors. Either an element was dropped (most likely
$\{4\}$ should read $\{4,6\}$, which would give $2+2+2=6$) or the printed
weight is for a different decomposition; and the listed sets contain $6>n_K=5$,
which the stated index range $1\le i\le n_K$ excludes. I cannot decode it from
the text I hold and therefore did not use it. This is the only datum in the
paper I wanted and could not obtain.
\end{remark}

\begin{remark}[numbering audit, and a self-inflicted scare]\label{rem:numbering}
Every HKKO number cited here was re-counted from the source this session rather
than recalled. HKKO declare \verb|\newtheorem{thm}{Theorem}[section]| with every
other environment sharing that counter, so numbers are determined by the order of
\verb|\begin{...}| lines. Counted: \texttt{defn:cylindric} is Defn.~2.1,
\texttt{def:par} Defn.~2.2, \texttt{prop:cssyt=crst} Prop.~2.4,
\texttt{defn:cSchur} Defn.~2.5, \texttt{prop:cyl\_sch} Prop.~2.6,
\texttt{prop:cssyt=path} Prop.~2.7, \texttt{thm:ssyt2} Thm.~2.8, and
\texttt{thm:C} --- the theorem Warnaar cites --- Thm.~3.3.

My first count made \texttt{thm:C} Theorem~3.2, which would have meant Warnaar's
``[HKKO25, Theorem 3.3]'' named a different theorem and that the whole $t=1$ gate of
\cite{Cq405} rested on the wrong result. The error was mine: my grep pattern listed
\texttt{thm,cor,conj,lem,assumption,prop,defn,quest,problem,rem} and omitted
\texttt{rems}, and there is a \verb|\begin{rems}| at source line 1113 which consumes
counter~3.2. \emph{An enumeration of an environment list is an instrument, and a
missing alternative in it shifts every subsequent number by one.} Recorded because
the first reading was alarming and wrong, and because two numbers in \cite{Cq405}
are wrong in the other direction: it cites Prop.~2.9 for \texttt{prop:cyl\_sch}
(it is 2.6) and Thm.~2.10 for \texttt{thm:ssyt2} (it is 2.8). Its Prop.~2.7 for the
path model is right. No statement in either paper depends on a number: every input
was quoted and, in the case of eq.~(3.6), verified symbolically at six parameter
points against three deliberately wrong variants.
\end{remark}

\begin{remark}[two of Korff's own standing restrictions exclude the target]
\label{rem:range}
Korff assumes $n_K>2$ throughout, which by Proposition~\ref{prop:dict} is
$2k>2$, i.e.\ $k\ge2$; the witness family of \eqref{eq:witness} has $k=1$, so it
sits at $n_K=2$. And Korff's Def.~5.5 defines
$P_{\la/d/\mu}(x_1,\dots,x_{\ell_K};t)$ for $1\le\ell_K\le k_K$ and sets it to
zero for $\ell_K>k_K$; here $\ell_K=n$ is the number of variables and
$k_K=w$, and the witness has $n=N=2k+w>w$, so Korff's own convention makes his
function vanish identically at every parameter where Warnaar's question has
content. I drop that convention --- formula \eqref{eq:Psi} is defined for any
number of variables --- and test the \emph{formula}, because answering a
mathematical question by a convention would be no answer. The combinatorics is
visibly non-empty there: e.g.\ $h=w=2$, $n=4$ has $4$ cylindric tableaux of
content $(1^4)$. These two restrictions are recorded, not used.
\end{remark}

\section{The obstruction: order of vanishing at $t=1$}

\begin{definition}\label{def:psitype}
Let $\nu(P):=\sum_{a=1}^{n}\bigl|J(\theta^{(a)})\bigr|$. A weight assignment
$\pi:\T\to\Z[t,t^{-1}]$ is \emph{cyclically $\psi$-graded} if for every $P$
either $\pi_P=0$ or $\ord_{t=1}\pi_P\ge\nu(P)$, where $\ord_{t=1}f$ denotes the
multiplicity of $t=1$ as a root of $f$.
\end{definition}

\begin{lemma}\label{lem:kor-is-psi}
For every $P$ and every $s(P)\in\Z$, $\varepsilon(P)\in\Z\setminus\{0\}$, the
weight $\pi_P=\varepsilon(P)\,t^{s(P)}\Psi_P(t)$ is cyclically $\psi$-graded,
with $\ord_{t=1}\pi_P=\nu(P)$ exactly.
\end{lemma}

\begin{proof}
By Lemma~\ref{lem:wd}, $\Psi_P$ is a product of exactly $\nu(P)$ factors
$(1-t^{m})$ with $m\ge1$, each vanishing to order exactly $1$ at $t=1$; and
$\varepsilon(P)t^{s(P)}$ is a unit at $t=1$.
\end{proof}

The next lemma is the whole mechanism, and the only place cyclicity is used.

\begin{lemma}[cyclic ascent]\label{lem:ascent}
For $\theta\in\{0,1\}^{h}$ one has $J(\theta)=\emptyset$ if and only if
$\theta$ is constant, i.e.\ $\theta=(0,\dots,0)$ or $\theta=(1,\dots,1)$.
Hence $\nu(P)\ge\#\{a:\theta^{(a)}\ \text{non-constant}\}
=\#\{a:\alpha_a\notin\{0,h\}\}$.
\end{lemma}

\begin{proof}
If $\theta$ is constant there is no $i$ with $\theta_i=0,\theta_{i+1}=1$.
Conversely let $\theta$ be non-constant, so the sets
$Z=\{i:\theta_i=0\}$ and $O=\{i:\theta_i=1\}$ are both non-empty. Pick
$p\in O$. Walking backwards from $p$ around the cycle $\Z/h\Z$ one must leave
$O$, since $Z\ne\emptyset$; let $i$ be the first index met with
$\theta_i=0$. Then $\theta_{i+1}=1$, so $i\in J(\theta)$. The last equality
holds because $\sum_i\theta_i=\alpha_a$, and a $0/1$ vector of length $h$ is
constant exactly when its sum is $0$ or $h$.
\end{proof}

\begin{lemma}\label{lem:nu}
Let $h\ge2$ and $\alpha=(1^{n})$. Then $\nu(P)=n$ for every $P\in\T_\alpha$.
\end{lemma}

\begin{proof}
Each step has $\sum_i\theta^{(a)}_i=1$, so $\theta^{(a)}$ has a single $1$, at
some position $p_a$; since $h\ge2$ it is non-constant. The only $i\in\Z/h\Z$
with $\theta^{(a)}_i=0$ and $\theta^{(a)}_{i+1}=1$ is $i=p_a-1$, so
$|J(\theta^{(a)})|=1$ for each of the $n$ steps.
\end{proof}

\begin{theorem}\label{thm:main}
Let $k\ge1$, $\ell\ge0$, $h=2k$, $w=2\ell+2$, $n\ge1$, and let
$\pi:\T\to\Z[t,t^{-1}]$ be cyclically $\psi$-graded. Suppose
\begin{equation}\label{eq:ansatz}
 (x_1\cdots x_n)^{k}\,\mathcal G_{k,\ell,n}(x;t)
 \;=\;\sum_{P\in\T}\pi_P(t)\,x^{\cont P}.
\end{equation}
Then for every content $\alpha$ with $c_\alpha\ne0$,
\[
  \ord_{t=1}c_\alpha(t)\ \ge\ \min\{\nu(P):P\in\T_\alpha,\ \pi_P\ne0\}.
\]
\end{theorem}

\begin{proof}
Comparing the coefficient of $x^{\alpha}$ in \eqref{eq:ansatz} gives
$c_\alpha=\sum_{P\in\T_\alpha}\pi_P$. Each summand is either $0$ or divisible by
$(t-1)^{\nu(P)}$ in $\Z[t,t^{-1}]$, hence the sum is divisible by
$(t-1)^{\min\nu(P)}$, the minimum over those $P$ with $\pi_P\ne0$. As
$c_\alpha\ne0$ its order of vanishing is finite and at least that exponent.
\end{proof}

\begin{corollary}[Korff's weight fails, uniformly in $\ell$]\label{cor:main}
Let $k=1$, $\ell\ge0$, $M=\ell+2$, $n=N=2M$, $\alpha=(1^{n})$. There is no
$\varepsilon:\T\to\Z$ and no $\mathrm{stat}:\T\to\Z$ with
\[
 (x_1\cdots x_n)\,\mathcal G_{1,\ell,n}(x;t)
 =\sum_{P\in\T}\varepsilon(P)\,t^{\mathrm{stat}(P)}\,\Psi_P(t)\,x^{\cont P}.
\]
In particular Huh--Kim--Krattenthaler--Okada's Theorem 3.3 with the cylindric
Schur function replaced by Korff's cylindric Hall--Littlewood function --- with
or without an additional monomial statistic, with HKKO's signs
\eqref{eq:cminus} or any others --- is false for every $\ell\ge0$.
\end{corollary}

\begin{proof}
Suppose such $\varepsilon,\mathrm{stat}$ exist. By
Lemma~\ref{lem:kor-is-psi} the weight is cyclically $\psi$-graded, and
$\T_\alpha\ne\emptyset$ (it contains the path with
$\theta^{(a)}=e_1$ for $a\le M$ and $e_2$ for $a>M$, whose endpoint is
$(M,M)$). By \eqref{eq:witness}, $c_\alpha\ne0$. By Lemma~\ref{lem:nu},
$\nu(P)=n=2M$ for every $P\in\T_\alpha$, since $h=2k=2\ge2$. So
Theorem~\ref{thm:main} forces
\begin{equation}\label{eq:need}
 \ord_{t=1}\bigl(C_M-(2M-1)t^{M-1}+t^{M+1}\bigr)\ \ge\ 2M .
\end{equation}
Evaluate the left side. At $t=1$ the polynomial takes the value
$C_M-2M+2$. For $M=2$ this is $2-4+2=0$, and indeed
$2-3t+t^{3}=(1-t)^{2}(t+2)$, whose order of vanishing at $t=1$ is exactly $2$
because $(t+2)|_{t=1}=3\ne0$; and $2<4=2M$. For $M\ge3$ we claim
$C_M>2M-2$, so that the order is $0<2M$. Indeed $C_3=5>4$, and
$C_{M+1}/C_M=2(2M+1)/(M+2)\ge2$ for $M\ge1$, so $C_M\ge5\cdot2^{M-3}$ for
$M\ge3$, while $2M-2<5\cdot2^{M-3}$ for all $M\ge3$ (at $M=3$: $4<5$; and the
right side at least doubles while the left grows by $2$). In either case
\eqref{eq:need} fails.
\end{proof}

\begin{remark}[the same argument at $k\ge2$]
Lemmas \ref{lem:ascent} and \ref{lem:nu} are uniform in $k$, so
Theorem~\ref{thm:main} applies verbatim; only the closed form
\eqref{eq:witness} is special to $k=1$. Computed values of $c_\alpha$ at
$\alpha=(1^n)$ from \eqref{eq:gen}:
$(k,\ell,n)=(2,0,4)$ gives $3t^{2}-6t+3=3(t-1)^2$, order $2$, against the
required $4$; $(2,1,6)$ gives $5t^{3}-15t^{2}+14$, order $0$, against the
required $6$. Two further rows are \textbf{vacuous and are recorded as such}:
at $(2,0,5)$ and $(2,1,7)$ one has $c_\alpha\equiv0$, so
Theorem~\ref{thm:main} says nothing there. Honest count: of six
$(k,\ell,n)$ tested with $k\ge2$, four are informative and two are empty.
\end{remark}

\section{Where the failure lives, exactly}

The obstruction is not that the target dislikes $\psi$-weights. At $\ell=0$ it
\emph{is} a $\psi$-weighted tableau sum --- over \emph{ordinary} tableaux.

\begin{proposition}\label{prop:located}
At $k=1$, $\ell=0$, $n=4$, $\alpha=(1^4)$:
\begin{enumerate}
\item Warnaar's Conjecture 1.5 asserts
$(x_1\cdots x_4)\mathcal G_{1,0,4}(x;t)=\sum_{\la\ \mathrm{even},\,\la_1\le2}P_\la(x;t)$,
whose $x^{\alpha}$-coefficient is $\sum_{T}\psi_T(t)$ over the two ordinary
semistandard tableaux of shape $(2,2)$ and content $(1^4)$. Their Macdonald
weights are $(1-t)^{2}$ and $(1-t)^{2}(1+t)$; each vanishes to order exactly
$2$ at $t=1$, i.e.\ at exactly $2$ of its $4$ horizontal strips; and their sum
is $(1-t)^{2}(t+2)=2-3t+t^{3}=c_\alpha(t)$.
\item The four $(2,2)$-cylindric tableaux of the same content all carry the
\emph{same} Korff weight $\Psi_P=(1-t)^{2}(1-t^{2})^{2}$, of order $4$ at
$t=1$: the cyclic ascent set fires at \emph{every} one of the $4$ strips.
\item Consequently $\sum_P c^{-}(\la(P))\Psi_P=0\ne c_\alpha(t)$, and no
statistic repairs it: $c_\alpha/\Psi=(t+2)/\bigl((1-t)^{2}(1+t)^{2}\bigr)$ is
not a Laurent polynomial.
\end{enumerate}
\end{proposition}

So the entire discrepancy is the integer $2$ versus $4$: passing from
Macdonald's open ascent set $\{1\le i<h:\theta_i=0,\theta_{i+1}=1\}$ to Korff's
cyclic one $\{i\in\Z/h\Z:\dots\}$ adds the wrap index $i=h$, and for a strip of
size $1$ in a $2$-row cylinder the wrap index fires exactly when the open one
does not. Each of the $4$ strips therefore contributes a vanishing factor
instead of $2$ of them doing so.

\begin{proposition}[the two deformations are normalised at opposite ends]
\label{prop:ends}
$\Psi_P(0)=1$ for every $P$ \textup{(Lemma~\ref{lem:wd})}, so Korff's
$P_{\la/d/\mu}(x;0)$ is the unweighted cylindric Schur function; whereas
\eqref{eq:C-} says that it is at $t=1$ that Warnaar's determinant becomes the
unweighted cylindric count. Accordingly
$\sum_{P\in\T_\alpha}c^{-}(\la(P))\Psi_P(0)=A_\alpha-B_\alpha=c_\alpha(1)$ for
every $\alpha$: verified at $(k,\ell,n)=(1,0,4)$ and $(1,1,6)$, where the two
sides are $0=0$ and $1=1$.
\end{proposition}

\begin{corollary}[what a correct weight must look like]\label{cor:survives}
Let $\pi$ be any weight satisfying \eqref{eq:ansatz}. At $k=1$,
$\alpha=(1^{2M})$ there must exist $P\in\T_\alpha$ with
$\ord_{t=1}\pi_P\le2$ ($M=2$) and with $\ord_{t=1}\pi_P=0$ ($M\ge3$). So for
$M\ge3$ at least one cylindric tableau of the fibre carries a weight that does
\emph{not} vanish at $t=1$, hence is not a product of factors $(1-t^{m})$,
$m\ge1$, at all. A weight of the surviving type must be a ratio with as many
such factors above as below --- a Gaussian-binomial type, whose value at $t=1$
is the corresponding ordinary binomial --- so that the $t=1$ specialisation
returns HKKO's unweighted count. This is a genuine narrowing: it rules out
$\psi_T$ and $\varphi_T$ and Korff's $\Psi_T,\Phi_T$ together with every
monomial regrading of them, and it points at the modified Hall--Littlewood /
Kostka--Foulkes normalisation rather than at the branching coefficients.
\end{corollary}

\section{Verification}

All computations are exact (SymPy, rational arithmetic); code in
\texttt{proofs/code-1010-q410/}.

\begin{enumerate}
\item \textbf{Ground truth re-derived.} $c_\alpha$ at $\alpha=(1^n)$ computed
directly from \eqref{eq:gen} by the determinant engine agrees with
\eqref{eq:witness} at $(k,\ell,n)=(1,0,4)$ and $(1,1,6)$:
$t^{3}-3t+2$ and $t^{4}-5t^{2}+5$, with $\ord_{t=1}=2$ and $0$.
\item \textbf{$\nu(P)=n$ at $\alpha=(1^n)$} (Lemma~\ref{lem:nu}) verified by
enumeration: the set of values of $\nu$ over all $P\in\T_{(1^n)}$ is the
singleton $\{n\}$ at $(k,\ell,n)=(1,0,4),(1,1,6),(1,2,8),(2,0,4),(2,1,6),(3,0,6)$,
over $4,18,68,6,64,20$ tableaux respectively.
\item \textbf{The $t=1$ evaluation, brute force.} Independently of
Lemma~\ref{lem:ascent}, for each content $\alpha$ the signed count
$\sum\{c^{-}(\la(P)):\Psi_P(1)\ne0\}$ was computed by evaluating
\eqref{eq:Psi} at $t=1$ on every tableau, and compared with the prediction
``$+1$ if every $\alpha_a\in\{0,h\}$, else $0$'' that
Lemma~\ref{lem:ascent} gives. Agreement in \textbf{23891 of 23891} contents at
$(k,\ell,n)=(1,0,4),(1,1,6),(1,0,3),(1,1,5),(2,0,4),(2,1,6),(1,2,8)$.
\item \textbf{The violations.} Comparing that signed count with
$c_\alpha(1)=A_\alpha-B_\alpha$:

\begin{center}
\begin{tabular}{lrrr}
\hline
$(k,\ell,n)$ & contents & contents violating & first violating $\alpha$ \\
\hline
$(1,0,3)$ & 27 & \textbf{0} & --- \\
$(1,0,4)$ & 81 & \textbf{0} & --- \\
$(1,1,5)$ & 243 & 90 & $(0,0,0,1,1)$, $c_\alpha(1)=1$ vs $0$ \\
$(1,1,6)$ & 729 & 301 & $(0,0,0,0,1,1)$, $c_\alpha(1)=1$ vs $0$ \\
$(1,2,8)$ & 6561 & 3025 & $(0^6,1,1)$, $c_\alpha(1)=1$ vs $0$ \\
$(2,0,4)$ & 625 & 65 & $(0,0,0,2)$, $c_\alpha(1)=1$ vs $0$ \\
$(2,1,6)$ & 15625 & 7749 & $(0^4,0,2)$, $c_\alpha(1)=1$ vs $0$ \\
\hline
\end{tabular}
\end{center}

The two $\ell=0$ rows are \textbf{zero}, and that is not a failure of the
method but a fact: at $h=w=2$ the walk $d=x_1-x_2$ alternates between
$\{0,2\}$ and $\{1\}$, so $A_\alpha=B_\alpha$ whenever any $\alpha_a=1$, and
the $t=1$ evaluation is satisfied identically. At $\ell=0$ the obstruction is
instead the order-$2$-versus-$4$ statement of
Proposition~\ref{prop:located}(3); for $\ell\ge1$ the $t=1$ evaluation alone
suffices. Both are instances of Theorem~\ref{thm:main}.
\item \textbf{Proposition~\ref{prop:q416a}} ($\Psi=\psi\cdot$ one wrap factor)
checked on \textbf{934 of 934} cylindric horizontal strips at
$(h,w)=(2,2),(2,4),(3,2),(3,3),(4,2),(4,4),(5,3)$, zero failures; with the wrap
factor omitted, $27$ strips disagree, so the check is not vacuous.
\item \textbf{Gates on the implementation of \eqref{eq:Psi}}, with their
discrimination controls, in \S\ref{sec:dict}: Korff (5.16) passed $180/180$
with all four wrong readings failing; non-cylindric reduction departing from
Macdonald's $\psi$ only at order $6,7$ with all wrong readings departing at
order $0$.
\end{enumerate}

\section{Q416, the wrap-factor conjecture: first half true, conclusion false}

This session's result also settles a conjecture opened a few hours earlier the same
day by my own browse slot, and it is worth separating what was already known from
what is proved here.

\paragraph{Already found, this morning, and not by me-now.} The browse slot of
2026-10-10 c3 recorded (i) that in Ram's and Schwer's alcove-walk form of the
Hall--Littlewood polynomials the $\psi/\varphi$ weight is $(1-t)^{\#\text{folds}}$,
one polynomial per positively folded walk, so a $\psi$-type weight \emph{is} a fold
count; and (ii) the $\ell=0$ capacity constraint that is the special case of
Corollary~\ref{cor:survives} it could see without a dictionary: $2-3t+t^{3}=(1-t)^2(t+2)$
has $t=1$ as a root of multiplicity \emph{exactly} $2$ (indeed $P''(1)=6$), and a
product of $r$ factors $(1-t^{m})$ vanishes to order exactly $r$, so every one of the
four cylindric tableaux must carry at least $2$ folds and at least one of them exactly
$2$. That is the same invariant as Theorem~\ref{thm:main}, found independently and
first, at one value of $\ell$ and as a constraint on a hypothetical statistic rather
than as a theorem about a published weight. It is recorded here as prior work of the
same day.

\paragraph{Q416, verbatim from the allocator.} ``\emph{$\Phi^{\mathrm{cyl}}_T(t)=
\psi_T(t)$ \textup{[Macdonald III (5.11$'$)]} times a product of $\psi$-type factors
indexed by the pairs where the wrap is \textsc{active} --- i.e.\ the polynomial weight
Q410 wants is the finite weight times a wrap factor, \textsc{not} a new statistic.}''

\begin{proposition}[the structural half is true]\label{prop:q416a}
Under the dictionary of Proposition~\ref{prop:dict}, Korff's weight is Macdonald's
weight times exactly one wrap factor per strip:
\[
 \Psi(x,\theta)=\psi(x,\theta)\cdot
 \begin{cases}
  \bigl(1-t^{\,w-(x_1-x_h)}\bigr), & \theta_h=0\ \text{and}\ \theta_1=1,\\
  1,&\text{otherwise},
 \end{cases}
\]
where $\psi(x,\theta)=\prod\{(1-t^{x_i-x_{i+1}}):1\le i<h,\ \theta_i=0,\theta_{i+1}=1\}$
is Macdonald's $\psi_{\la/\mu}$ for the corresponding ordinary strip.
\end{proposition}

\begin{proof}
Korff's index set is $J(\theta)=\{i\in\Z/h\Z:\theta_i=0,\theta_{i+1}=1\}$ and
Macdonald's is the same condition for $1\le i<h$ with $\theta_{h+1}:=0$, which can
never put $i=h$ in the set. The two differ exactly in the single index $i=h$, which
lies in $J(\theta)$ iff $\theta_h=0$ and $\theta_{h+1}=\theta_1=1$, and its factor is
$(1-t^{m_h(x)})=(1-t^{\,w-(x_1-x_h)})$ by \eqref{eq:mvec}. The wrap index is active
precisely when the strip wraps, i.e.\ when the step increases $x_1$ without increasing
$x_h$.
\end{proof}

Measured consistency: the two weights therefore agree to order $w-\mathrm{span}$ in
$t$, and this is Gate 2 of \S\ref{sec:dict} --- departure orders $6$ and $7$ at
$w=8,9$, against order $0$ for every misplaced wrap.

\begin{corollary}[the conclusion is false]\label{cor:q416b}
Q416's conclusion fails. By Lemma~\ref{lem:nu} the wrap factor is active at
\emph{every} strip of \emph{every} cylindric tableau of content $(1^{n})$: at
$\ell=0$, $k=1$ each of the four tableaux has fold count $4$, not $2$, and all four
carry the same weight $(1-t)^{2}(1-t^{2})^{2}$. So Korff's weight violates the
browse slot's own $\ell=0$ constraint --- which demanded at least one tableau with
fold count exactly $2$ --- and it does so because the wrap factor pushes the order of
vanishing in the \emph{wrong direction}: from the $2$ that Macdonald's $\psi$ supplies
on the two ordinary tableaux of shape $(2,2)$ \textup{(Proposition~\ref{prop:located})}
up to $4$. ``The finite weight times a wrap factor'' is exactly the object that cannot
work.
\end{corollary}

This is the useful form of the negative result: a correct weight must make the
cylindric wrap \emph{not} count as a fold, or must cancel it. The fold count is the
quantity to control, and it must be bounded by $2$ on some tableau of each fibre, not
driven up by the periodicity.

\section{Status of Q410}

\textbf{Answered: no}, and Q416 with it (Proposition~\ref{prop:q416a},
Corollary~\ref{cor:q416b}: its structural half is a theorem, its conclusion is
refuted). Korff's cylindric Hall--Littlewood weight is not the
polynomial Warnaar's \S6 needs, and nothing cyclically $\psi$-graded is. The
gap between ``cylindric Hall--Littlewood functions exist'' and ``Warnaar's
determinant is a weighted cylindric tableau sum'' is therefore real, and it is
now located at a single invariant: the order of vanishing at $t=1$, which the
cylindric ascent set forces up to the number of non-constant strips while the
target vanishes to order at most $2$.

Two by-products. (i) Proposition~\ref{prop:ends}: Korff's deformation of the
cylindric Schur function and Warnaar's deformation of the same object are
normalised at opposite values of $t$, and $\sum_P c^-\Psi_P(0)=c_\alpha(1)$
is the precise statement of that. (ii) Corollary~\ref{cor:survives} names the
surviving type. The scoop risk recorded in the brief --- that HKKO own both
halves of Warnaar's \S6 and their Lemma 7.4 connects them --- is unaffected by
a negative result, but note that neither HKKO nor Warnaar cites Korff
(measured: ``Korff'' occurs $0$ times in HKKO's 41 bibitems and $0$ times in
Warnaar's 183{,}950 characters), which the present note explains rather than
laments.

\begin{thebibliography}{9}
\bibitem{HKKO} J.~Huh, S.~Kim, C.~Krattenthaler, S.~Okada,
\emph{Two affine bounded Littlewood identities: symplectic case},
arXiv:2301.13117. (Full \LaTeX{} source held;
Thm.~3.3 (\texttt{thm:C}), eq.~(3.6) (\texttt{eq:C-}),
Defn.~2.1, Defn.~2.2, Prop.~2.4, Defn.~2.5, Prop.~2.6, Prop.~2.7, Thm.~2.8
read verbatim.  Numbers re-counted this session, see Remark~\ref{rem:numbering}.)
\bibitem{K} C.~Korff, \emph{Cylindric versions of specialised Macdonald
functions and a deformed Verlinde algebra}, arXiv:1110.6356v3, 6 Sep 2012.
(Full PDF held, 77 pp.; \S\S2.19--2.22 and 5.2--5.3, Defs.~5.1--5.5,
eqs.~(5.7)--(5.17), Lemmas 5.4--5.5 read verbatim.)
\bibitem{M} I.~G.~Macdonald, \emph{Symmetric Functions and Hall Polynomials},
2nd ed., Oxford. (Held; III (5.8$'$), (5.9$'$), (5.11$'$), pp.~229--230.)
\bibitem{W} S.~O.~Warnaar, arXiv:2511.17034 (SIGMA 26 (2026), 062).
(Full source held; \S6 and Conjecture 1.5 read verbatim.)
\bibitem{Cq405} Clio, \emph{No monomial statistic on cylindric tableaux carries
Warnaar's generalised affine Jacobi--Trudi determinant},
\texttt{proofs/2026-10-10-Q405-twisted-trace-cylindric-statistic.tex},
10 October 2026.
\end{thebibliography}

\end{document}
